\documentclass[12pt]{article}
\usepackage[utf8]{inputenc}
\usepackage[T1]{fontenc}
\usepackage[margin=1in]{geometry}
\usepackage{setspace}
\usepackage{xcolor}
\usepackage[colorlinks=true, urlcolor=blue!60!black, linkcolor=blue!60!black]{hyperref}
\setlength{\parindent}{2em}
\setlength{\parskip}{1.25em}
\renewcommand{\baselinestretch}{1.0}
\begin{document}
\begin{center}
{\LARGE \textit{Hernandez v. Texas}, 347 U.S. 475 (1954)\par}
\vspace{0.8em}
{\large Case Brief\par}
\end{center}
\vspace{1em}

\section{Memory Jogger}
The Fourteenth Amendment's Equal Protection Clause is not limited to a Black--white ``two-class theory,'' and bars the systematic exclusion of Mexican Americans---or any other locally subordinated distinct class---from jury service (\href{https://www.courtlistener.com/opinion/105220/hernandez-v-texas/}{CourtListener}).

\section{Detailed Case Facts}
In 1951, petitioner Pete Hernandez, a Mexican-American agricultural worker, was indicted for the murder of Joe Espinosa by an all-Anglo grand jury in Jackson County, Texas, convicted by an all-Anglo petit jury, and sentenced to life imprisonment (\href{https://www.oyez.org/cases/1940-1955/347us475}{Oyez}; \href{https://supreme.justia.com/cases/federal/us/347/475/}{Justia}). Represented \textit{pro bono} by a team of Mexican-American civil rights attorneys---including Gustavo C.\ (``Gus'') Garc\'{i}a, Carlos C.\ Cadena, John J.\ Herrera, and James DeAnda, making this the first U.S.\ Supreme Court case argued by Mexican-American lawyers---Hernandez challenged the systematic exclusion of persons of Mexican descent from jury commissions, grand juries, and petit juries in Jackson County (\href{https://www.oyez.org/cases/1940-1955/347us475}{Oyez}; \href{https://scholar.smu.edu/cgi/viewcontent.cgi?article=1879&context=law_faculty}{SMU Law Review}).

According to the 1950 U.S.\ Census, 1,865 of Jackson County's 12,916 residents (approximately 14\%) bore Mexican or Latin American surnames, including 1,738 native-born U.S.\ citizens and 65 naturalized citizens; among the county's 3,754 males over age 21, 408 (about 11\%) had Spanish surnames (\href{https://supreme.justia.com/cases/federal/us/347/475/}{Justia}). The County Tax Assessor testified that 6\% to 7\% of the freeholders on the county tax rolls were of Mexican descent (\href{https://supreme.justia.com/cases/federal/us/347/475/}{Justia}). Crucially, the State of Texas stipulated to two facts: (1) for the preceding 25 years, there was no record of any person with a Mexican or Latin American name serving on a jury commission, grand jury, or petit jury in Jackson County (out of more than 6,000 jurors called), and (2) there were male persons of Mexican or Latin American descent in Jackson County who met every legal prerequisite---including citizenship and freeholder status---to serve (\href{https://supreme.justia.com/cases/federal/us/347/475/}{Justia}; \href{https://caselaw.findlaw.com/court/us-supreme-court/347/475.html}{FindLaw}).

Although Texas's statutory jury-commission system (Vernon's Tex.\ Code Crim.\ Proc.\ arts.\ 333--350; Tex.\ Civ.\ Stat.\ art.\ 2107) was facially neutral, petitioner presented extensive evidence of local social segregation demonstrating that persons of Mexican descent were treated as a distinct class apart from ``whites'' in Jackson County (\href{https://supreme.justia.com/cases/federal/us/347/475/}{Justia}). Officials and residents admitted the community distinguished between ``white'' and ``Mexican''; participation by Mexican Americans in local business and civic groups was slight; until very recently, children of Mexican descent were required to attend a segregated school for the first four grades; at least one restaurant in town prominently displayed a sign reading ``No Mexicans Served''; and on the Jackson County courthouse grounds at the time of the hearing, there were two men's restrooms---one unmarked, and the other marked ``Colored Men'' and ``Hombres Aqui'' (``Men Here'') (\href{https://supreme.justia.com/cases/federal/us/347/475/}{Justia}). To rebut petitioner's showing, the State offered only the testimony of five jury commissioners who asserted that they had not discriminated and had sought solely to pick the best-qualified jurors (\href{https://supreme.justia.com/cases/federal/us/347/475/}{Justia}).

\section{Procedural History}
Before trial in the District Court of Jackson County, Texas, Hernandez's counsel filed timely motions to quash both the indictment and the petit jury panel, alleging that the systematic exclusion of persons of Mexican descent violated the Equal Protection Clause of the Fourteenth Amendment (\href{https://supreme.justia.com/cases/federal/us/347/475/}{Justia}). After an evidentiary hearing, the trial court denied the motions; counsel renewed them at trial with additional evidence, and the court again denied them (\href{https://supreme.justia.com/cases/federal/us/347/475/}{Justia}). Hernandez was convicted of murder and sentenced to life imprisonment (\href{https://www.oyez.org/cases/1940-1955/347us475}{Oyez}).

On appeal---resting solely on the denial of the jury-exclusion motions---the Texas Court of Criminal Appeals affirmed (\textit{Hernandez v.\ State}, 160 Tex.\ Crim.\ 72, 251 S.W.2d 531 (1952)) (\href{https://supreme.justia.com/cases/federal/us/347/475/}{Justia}). Relying on prior Texas precedent, the state appellate court held that the Fourteenth Amendment recognized only two classes (``white'' and ``Negro'') and that because Mexican Americans were classified as members of the white race, their exclusion from juries did not violate equal protection (\href{https://www.oyez.org/cases/1940-1955/347us475}{Oyez}; \href{https://lawcat.berkeley.edu/record/1116057/files/fulltext.pdf}{Berkeley Law Scholarship}). The U.S.\ Supreme Court granted certiorari (346 U.S.\ 811 (1953)), heard oral argument on January 11, 1954, and reversed unanimously on May 3, 1954---just two weeks before \textit{Brown v.\ Board of Education}, 347 U.S.\ 483 (1954) (\href{https://supreme.justia.com/cases/federal/us/347/475/}{Justia}). (On retrial before a jury that included Mexican Americans, Hernandez was convicted again (\href{https://cl.cobar.org/departments/latinos-in-the-courtroom/}{Colorado Lawyer}).)

\section{Judicial Votes}
\begin{itemize}
\item \textbf{Judgment:} Unanimous (9--0) reversal of the Texas Court of Criminal Appeals (\href{https://www.oyez.org/cases/1940-1955/347us475}{Oyez}).
\item \textbf{Opinion of the Court (All 9 Justices):} Chief Justice Earl Warren delivered the opinion of the Court, joined by Justices Hugo L.\ Black, Stanley F.\ Reed, Felix Frankfurter, William O.\ Douglas, Robert H.\ Jackson, Harold H.\ Burton, Tom C.\ Clark, and Sherman Minton (\href{https://www.oyez.org/cases/1940-1955/347us475}{Oyez}).
\item \textbf{Concurrences:} None (\href{https://supreme.justia.com/cases/federal/us/347/475/}{Justia}).
\item \textbf{Dissents:} None (\href{https://supreme.justia.com/cases/federal/us/347/475/}{Justia}).
\end{itemize}

\section{Holding}
\textbf{Judgment:} Reversed. The murder conviction was overturned because the systematic exclusion of persons of Mexican descent from jury service in Jackson County deprived Hernandez of the equal protection of the laws (\href{https://supreme.justia.com/cases/federal/us/347/475/}{Justia}).

\textbf{Rule Announced:} The Equal Protection Clause of the Fourteenth Amendment is not confined to a ``two-class theory'' of ``white'' and ``Negro''; it forbids the exclusion of otherwise eligible persons from jury service solely because of their ancestry or national origin (\href{https://supreme.justia.com/cases/federal/us/347/475/}{Justia}). To establish an equal protection violation for a non-Black/white group, a challenger must satisfy a two-step burden: (1) prove as a question of fact that the group constitutes a distinct, identifiable class within the local community, and (2) establish a \textit{prima facie} case of systematic exclusion under the ``rule of exclusion'' from \textit{Norris v.\ Alabama}, 294 U.S.\ 587 (1935) (substantial population presence, qualified members, and complete exclusion over an extended period) (\href{https://supreme.justia.com/cases/federal/us/347/475/}{Justia}). Once established, mere general denials of discriminatory intent by jury commissioners are insufficient to rebut the \textit{prima facie} case (\href{https://caselaw.findlaw.com/court/us-supreme-court/347/475.html}{FindLaw}).

\textbf{Holding vs.\ Dicta:} The Court held that persons of Mexican descent were proven to be a separate class in Jackson County and that their 25-year total exclusion among over 6,000 jurors established unconstitutional discrimination (\href{https://caselaw.findlaw.com/court/us-supreme-court/347/475.html}{FindLaw}). In footnote 9, Chief Justice Warren expressly reserved as undecided whether courts could take judicial notice that persons of Mexican descent constitute a separate class without local factual proof (\href{https://supreme.justia.com/cases/federal/us/347/475/}{Justia}). Additionally, the Court's reassurance that the Fourteenth Amendment does not require ``proportional representation of all the component ethnic groups of the community on every jury'' clarifies the outer limit of the right rather than resolving a claim actually pressed by Hernandez (\href{https://caselaw.findlaw.com/court/us-supreme-court/347/475.html}{FindLaw}).

\section{Analysis of Opinions}

\subsection{Majority Opinion (Chief Justice Warren)}
Chief Justice Warren dismantled Texas's ``two-class theory'' by situating the Fourteenth Amendment in both precedent and social reality. Citing \textit{Strauder v.\ West Virginia}, 100 U.S.\ 303 (1880) (noting that excluding naturalized ``Celtic Irishmen'' from juries would violate the Amendment's spirit), \textit{Truax v.\ Raich}, \textit{Takahashi v.\ Fish \& Game Commission}, and \textit{Hirabayashi v.\ United States}, Warren observed that ``community prejudices are not static, and, from time to time, other differences from the community norm may define other groups which need the same protection'' (\href{https://supreme.justia.com/cases/federal/us/347/475/}{Justia}). He also noted that Texas's own courts had recognized broader equal protection claims outside the Mexican-American jury context, including barring the exclusion of Roman Catholics from juries and refusing to enforce restrictive covenants against persons of Mexican descent (\href{https://supreme.justia.com/cases/federal/us/347/475/}{Justia}).

Turning to the two-part factual inquiry, Warren first held that Hernandez proved Mexican Americans were treated as a distinct class in Jackson County based on ``the attitude of the community''---specifically local testimony, segregated schooling, the ``No Mexicans Served'' restaurant sign, and the courthouse restroom marked ``Colored Men'' and ``Hombres Aqui'' (\href{https://supreme.justia.com/cases/federal/us/347/475/}{Justia}). Rejecting Texas's objection to using Spanish surnames as a proxy, Warren noted that jury commissioners worked from lists of names, so Spanish surnames provided ``ready identification'' just as skin color did for race (\href{https://supreme.justia.com/cases/federal/us/347/475/}{Justia}).

Second, Warren applied the \textit{Norris v.\ Alabama} ``rule of exclusion.'' Given that 14\% of the county population and 11\% of adult males had Spanish surnames, that 6\% to 7\% of freeholders were of Mexican descent, and that Texas stipulated qualified Mexican Americans existed yet zero had served among over 6,000 jurors in 25 years, ``it taxes our credulity to say that mere chance'' explained the outcome (\href{https://caselaw.findlaw.com/court/us-supreme-court/347/475.html}{FindLaw}). Under \textit{Norris}, the five jury commissioners' general assertions that they did not discriminate could not overcome this \textit{prima facie} proof, ``whether or not it was a conscious decision on the part of any individual jury commissioner'' (\href{https://caselaw.findlaw.com/court/us-supreme-court/347/475.html}{FindLaw}).

\subsection{Concurrences and Dissents}
No Justice filed a concurring or dissenting opinion; Chief Justice Warren spoke for a completely unanimous Court (\href{https://www.oyez.org/cases/1940-1955/347us475}{Oyez}).

\section{Examples: Future Applications}
\begin{enumerate}
\item \textbf{Same-Side (Equal Protection Violation):} In a county where Hmong-American citizens make up 12\% of the adult population and face documented local segregation in housing and civic life, jury commissioners select jurors from voter lists by hand, and zero Hmong-surnamed citizens have been called among 4,500 jurors over 20 years despite stipulated qualified residents. \textit{Analysis:} Directly tracks \textit{Hernandez}: local community attitudes establish a distinct class defined by ancestry/national origin, and prolonged total exclusion triggers the \textit{Norris} \textit{prima facie} inference that general denials cannot rebut (\href{https://supreme.justia.com/cases/federal/us/347/475/}{Justia}).
\item \textbf{Same-Side (Equal Protection Violation):} A state county with a substantial, locally ostracized Arab-American population uses a facially neutral ``key-man'' jury-selection system; over 15 years, not a single person of Arab descent is selected for a grand jury, and commissioners testify only that they picked ``respectable acquaintances.'' \textit{Analysis:} Under \textit{Hernandez}, a facially neutral selection statute is unconstitutional as applied when administered to systematically exclude a distinct national-origin class, regardless of whether commissioners claim benign motives (\href{https://caselaw.findlaw.com/court/us-supreme-court/347/475.html}{FindLaw}).
\item \textbf{Opposite-Side (No Equal Protection Violation):} A Mexican-American defendant is tried in a county where Mexican Americans comprise 18\% of the population and routinely serve on grand and petit juries in roughly proportional numbers, but by random draw none happens to sit on the defendant's 12-person petit jury. \textit{Analysis:} Chief Justice Warren explicitly held that \textit{Hernandez} does not guarantee proportional representation or a right to have members of one's group on a particular jury---only freedom from systematic exclusion (\href{https://caselaw.findlaw.com/court/us-supreme-court/347/475.html}{FindLaw}).
\item \textbf{Opposite-Side (No Equal Protection Violation):} A state excludes convicted felons and non-English speakers from jury service through uniform, facially neutral statutory qualifications, and no evidence shows that any ethnic group is singled out or systematically excluded beyond those neutral criteria. \textit{Analysis:} \textit{Hernandez} forbids differential treatment ``not based on some reasonable classification'' and applies only to ``otherwise eligible persons'' excluded because of ancestry or national origin (\href{https://supreme.justia.com/cases/federal/us/347/475/}{Justia}).
\item \textbf{Fence-Sitter (Borderline):} In a county where 30\% of residents are of Mexican descent, jury commissioners over a 10-year period select Mexican Americans to comprise 4\% of grand jurors (never zero, unlike Jackson County's 25-year total exclusion), while offering no explanation for the sharp statistical disparity. \textit{Analysis:} Unlike \textit{Hernandez}'s ``inexorable zero'' under \textit{Norris}, some class members did serve, creating a closer evidentiary dispute over whether severe underrepresentation suffices for a \textit{prima facie} case---a question the Supreme Court later resolved in favor of challengers using statistical disparity analysis in \textit{Castaneda v.\ Partida}, 430 U.S.\ 482 (1977) (\href{https://lawcat.berkeley.edu/record/1116303/files/fulltext.pdf}{Berkeley Law Scholarship}).
\end{enumerate}

\section{Critique}
\textbf{Originalist and Textualist Perspectives:} From a strict historical-intent perspective rooted in early Reconstruction decisions like the \href{https://fedsoc.org/case/the-slaughter-house-cases}{\textit{Slaughter-House Cases}, 83 U.S.\ (16 Wall.) 36 (1873)}, critics could observe that the Supreme Court initially doubted whether the Equal Protection Clause extended beyond discrimination against formerly enslaved African Americans (\href{https://law.justia.com/constitution/us/amendment-14/06-equal-protection-of-the-laws.html}{Justia Constitution Annotated}). Yet textualists broadly endorse \textit{Hernandez}'s rejection of the ``two-class theory'': the text of the Fourteenth Amendment protects ``any person'' within a state's jurisdiction, not merely two racial categories (\href{https://law.justia.com/constitution/us/amendment-14/06-equal-protection-of-the-laws.html}{Justia Constitution Annotated}). At the same time, formalist and colorblindness-oriented critics question Warren's open-ended sociological test (``community prejudices are not static'') and his willingness to find discrimination ``whether or not it was a conscious decision,'' which stands in tension with later Equal Protection doctrine requiring proof of discriminatory intent under \textit{Washington v.\ Davis} (1976) (\href{https://law.justia.com/constitution/us/amendment-14/06-equal-protection-of-the-laws.html}{Justia Constitution Annotated}).

\textbf{Progressive, LatCrit, and Doctrinal Perspectives:} Critical race and Latino legal scholars---most prominently Ian F.\ Haney L\'{o}pez, Kevin R.\ Johnson, and Michael A.\ Olivas---offer a nuanced, dual assessment of \textit{Hernandez} (\href{https://www.law.uh.edu/hernandez50/KevinJohnson.pdf}{Johnson, UH Law Center}; \href{https://lawcat.berkeley.edu/record/1116303/files/fulltext.pdf}{Haney L\'{o}pez, 85 Calif.\ L.\ Rev.\ 1143 (1997)}). On one hand, Haney L\'{o}pez praised \textit{Hernandez} as ``the single most insightful Supreme Court opinion on race ever handed down'' because it treated group identity and subordination as a dynamic social construct rather than a fixed biological essence (\href{https://scholarsarchive.byu.edu/cgi/viewcontent.cgi?article=1129&context=byuplr}{BYU Prelaw Review}; \href{https://escholarship.org/content/qt43t5940z/qt43t5940z.pdf}{Haney L\'{o}pez, 25 Chicano-Latino L.\ Rev.\ 61 (2005)}). On the other hand, scholars critique the decision's reliance on the ``other white'' / ``class apart'' litigation strategy: by refusing to declare Mexican Americans a protected racial group nationwide and instead requiring plaintiffs to prove local ``community attitudes'' county by county, \textit{Hernandez} imposed an onerous, repetitive evidentiary burden on Latino litigants and left the Black--white racial binary largely intact (\href{https://digitalcommons.law.byu.edu/cgi/viewcontent.cgi?article=1231&context=elj}{BYU Educ.\ \& L.J.}; \href{https://lawcat.berkeley.edu/record/1116303/files/fulltext.pdf}{Haney L\'{o}pez, 85 Calif.\ L.\ Rev.\ 1143}).

\textbf{Underlying Competing Values:} \textit{Hernandez} highlights three foundational tensions in Equal Protection jurisprudence: (1) anti-subordination (remedying real-world social hierarchy and local stigma) versus formal anti-classification (treating race as a fixed, symmetrical category); (2) federal constitutional supervision of state criminal justice systems versus state discretion over jury selection; and (3) statistical and structural proof of exclusion versus subjective, individual proof of discriminatory intent (\href{https://escholarship.org/content/qt43t5940z/qt43t5940z.pdf}{Haney L\'{o}pez, 25 Chicano-Latino L.\ Rev.\ 61}).

\section{Key Quotations}
\begin{quote}
``Throughout our history, differences in race and color have defined easily identifiable groups which have at times required the aid of the courts in securing equal treatment under the laws. But community prejudices are not static, and, from time to time, other differences from the community norm may define other groups which need the same protection.''
\end{quote}
\noindent --- Chief Justice Earl Warren, Opinion of the Court, 347 U.S.\ at 478 (\href{https://supreme.justia.com/cases/federal/us/347/475/}{Justia}).

\begin{quote}
``When the existence of a distinct class is demonstrated, and it is further shown that the laws, as written or as applied, single out that class for different treatment not based on some reasonable classification, the guarantees of the Constitution have been violated. The Fourteenth Amendment is not directed solely against discrimination due to a `two-class theory'---that is, based upon differences between `white' and Negro.''
\end{quote}
\noindent --- Chief Justice Earl Warren, Opinion of the Court, 347 U.S.\ at 478 (\href{https://supreme.justia.com/cases/federal/us/347/475/}{Justia}).

\begin{quote}
``On the courthouse grounds at the time of the hearing, there were two men's toilets, one unmarked, and the other marked `Colored Men' and `Hombres Aqui' (`Men Here').''
\end{quote}
\noindent --- Chief Justice Earl Warren, Opinion of the Court, 347 U.S.\ at 479--480 (\href{https://supreme.justia.com/cases/federal/us/347/475/}{Justia}).

\begin{quote}
``But it taxes our credulity to say that mere chance resulted in there being no members of this class among the over six thousand jurors called in the past 25 years. The result bespeaks discrimination, whether or not it was a conscious decision on the part of any individual jury commissioner.''
\end{quote}
\noindent --- Chief Justice Earl Warren, Opinion of the Court, 347 U.S.\ at 482 (\href{https://caselaw.findlaw.com/court/us-supreme-court/347/475.html}{FindLaw}).

\begin{quote}
``His only claim is the right to be indicted and tried by juries from which all members of his class are not systematically excluded---juries selected from among all qualified persons regardless of national origin or descent. To this much he is entitled by the Constitution.''
\end{quote}
\noindent --- Chief Justice Earl Warren, Opinion of the Court, 347 U.S.\ at 482 (\href{https://supreme.justia.com/cases/federal/us/347/475/}{Justia}).

\end{document}
